\bibitem{hodge-fargues-finite}
L.~Fargues, with the collaboration of Y.~Tian,
\emph{La filtration canonique des points de torsion des groupes
$p$-divisibles}, Annales scientifiques de l'\'Ecole normale
sup\'erieure \textbf{44} (2011), 905--961.
The base-field convention is on p.~911; the finite-group
Hodge--Tate theorem is \S5.4, Theorem~3, p.~935.
\url{https://doi.org/10.24033/asens.2157}.

\bibitem{hodge-finite-part}
The Stacks Project, \emph{Henselian local rings},
Lemma~10.153.5 (tag 04GJ), and \emph{Zariski's Main Theorem},
Lemma~37.43.3 (tag 05K0). Factoring through a finite scheme and
using its henselian decomposition gives the finite part of a
separated quasi-finite scheme over a henselian local base.
\url{https://stacks.math.columbia.edu/tag/04GJ};
\url{https://stacks.math.columbia.edu/tag/05K0}.

\bibitem{hodge-flattening}
The Stacks Project, \emph{Blowing up and flatness},
Lemma~31.36.3 (tag 0ESN).
\url{https://stacks.math.columbia.edu/tag/0ESN}.

\bibitem{hodge-specialization-source}
L.~Alp\"oge, original Hodge--Tate cotangent modules, finite torsion
orders and corrected-fibre comparison, \emph{jsmorph/rh}, commit
\texttt{ecd09ff9b4aaf6c08f2b5ef17dc49ca9648a77de}.
The exact declaration ranges and authenticated source hashes are
recorded in \texttt{certificates/geometry-hodge-specialization.json};
the source dictionary and the scope of the finite-group argument
are in \texttt{notes/geometry-hodge-specialization.md}.
\url{https://github.com/jsmorph/rh/tree/ecd09ff9b4aaf6c08f2b5ef17dc49ca9648a77de/lean/src}.
